\cotefile{91}
\scanpage{73}
% manuscript line 1
recollement en un morphisme est immédiat.

% manuscript line 2
Par suite, on a défini $W(X)\xrightarrow{\varphi}\unsure{\mathfrak{X}|\mathscr{X}}$.

% manuscript line 3
Je dis que c'est une solution du Pb universel .... i.e\unsure{,|.} que pour tt $T$,
% manuscript line 5
\[
\operatorname{Hom}_S(T,X)\longrightarrow
\operatorname{Hom}_{\underline W(S)}(\underline W(T),\mathfrak{X})
\]
% manuscript line 6
est \uline{bijectif}. Soit \deleted{$u:T\to X$},
% manuscript line 7
\[
\underline W(T)\xrightarrow{\unsure{v|\psi}}\mathfrak{X}
\]
% manuscript line 8
je dis qu'il peut se factoriser de \uline{façon unique} en
% manuscript line 10
\[
\underline W(T)\xrightarrow{\underline W(u)}\underline W(X)
\overset{\varphi}{\unsure{\text{---}|\longrightarrow}}\mathfrak{X}
\]

% manuscript line 11
En effet, \deleted{\unsure{considérons|supposons} \unsure{les|que} \illegible{1}} \deleted{il \unsure{faudr.|faudra} \unsure{que} \unsure{si|l'on ait} on peut factoriser, \unsure{ce} \unsure{soit} $T=\bigcup$} \deleted{$W(u)($} \deleted{\unsure{inversement}} utilisons la condition $(\beta)$\,: $T=\bigcup T_i$, avec $v(\underline W(T_i))\subset\mathfrak{X}_i$. Donc $v|\underline W(T_i))$ est défini par $\unsure{v_i|u_i}:T_i\to X_i$, par factorisation
% manuscript line 18
\[
\underline W(T_i)\xrightarrow{\underline W(u_i)}
 \underline W(X_i)\xrightarrow{\varphi_i}X_i
\]
% manuscript line 19
On vérifie aussitôt que les composés
% manuscript line 20
\[
 T_i\longrightarrow X_i\longrightarrow X
\]
\deleted{$\underline W(T_i)\to\underline W(X_i)\to\underline W(X)$}
% archive sheet number: 72
